%% kapa: exam=1 round=21 year=2010 period=2 form=A subject=accounting subject_ko=회계학 seq=22 number=22
%% source: Q-Net 감정평가사 기출문제 (cst00309) artlSeq=5024181 file=extracted/1008/5024181/2010년 제21회 감정평가사 1차시험 기출문제/2교시.hwp
%% conversion: kordoc 4.24.51 -> tools/gen1.py (auto); review=proofread-gpt-sol; flags=none
\begin{question}{22}
(주)감평은 20×0년 7월 말 당좌예금잔액을 은행에 조회한 결과, ￦4,325,000이라는 회신을 받았다. 한편, (주)감평의 정확한 당좌예금 잔액을 산출하기 위한 관련 자료는 아래와 같다. 아래에 대한 차이 조정 전 20×0년 7월 말 (주)감평의 장부상 당좌예금잔액은 얼마인가?
\begin{framedtext}
ㆍ(주)감평이 20×0년 7월 10일에 입금한 ￦760,000이 입금 처리되지 않았다. \par
ㆍ(주)감평이 20×0년 7월 15일에 발행한 수표 ￦830,000이 아직 인출되지 않았다. \par
ㆍ은행이 20×0년 7월 20일에 (주)국세의 당좌거래 수수료 ￦15,000을 잘못하여 (주)감평에 부과하고, 이를 (주)감평의 당좌예금에서 차감하였다. \par
ㆍ은행이 20×0년 7월 25일에 (주)공인에 귀속시켜야 할 이자수익 ￦30,000을 잘못하여 (주)감평의 당좌예금에 입금처리하였다. \par
ㆍ20×0년 7월 말 은행이 보내온 예금증명서에 포함되어 있는 금액 중 (주)감평의 장부에 반영되지 않은 것은 다음과 같다. \par
- 추심의뢰한 받을어음 ￦450,000이 (주)감평의 당좌예금에 입금됨 \par
- (주)감평의 거래처가 (주)감평에 대한 외상매입을 갚기 위하여 ￦340,000을 은행에 직접 송금함 \par
- (주)감평이 거래처로부터 받은 수표 ￦127,000을 당좌예금에 입금하였으나, 은행에서는 이를 최종부도 처리함 \par
\end{framedtext}
\begin{choices}
  \item ￦3,450,000
  \item ￦3,517,000
  \item ￦3,577,000
  \item ￦4,180,000
  \item ￦4,240,000
\end{choices}
\end{question}
\answer{3}
